% Diagrama: Arquitectura Ports-and-Adapters de VFTModel
% Usado en: 4-Capas-Codigo/4-2-VFTModel.tex  (fig:VFTModel-arquitectura)
% Fix 2026-08-28: eliminados actores externos como cajas separadas (diagrama era >18cm);
%   información de sistemas externos integrada dentro de los adaptadores.
\begin{tikzpicture}[
    dominio/.style = {
        rectangle, draw=unamAzul, fill=unamAzul!18,
        text width=3.8cm, align=center,
        minimum height=2.4cm, rounded corners=5pt,
        font=\small\bfseries
    },
    adaptador/.style = {
        rectangle, draw=unamAzul!70, fill=unamAzul!6,
        text width=3.2cm, align=center,
        minimum height=1.8cm, rounded corners=3pt,
        font=\small
    },
    flecha/.style = {Stealth-Stealth, thick, color=unamAzul!55}
]

%% — Dominio central —
\node[dominio] (dom) {
    \textbf{Dominio}\\
    \texttt{src/core/}\\[3pt]
    {\footnotesize Algoritmos \textsc{vft}\\Constructor del grafo}
};

%% — Adaptadores alrededor del dominio —
\node[adaptador, left=0.8cm of dom] (adp-api) {
    \textbf{Adaptador \api{}}\\
    \texttt{src/api/}\\[3pt]
    {\footnotesize FastAPI \textsc{rest} $\cdot$ Go}
};
\node[adaptador, right=0.8cm of dom] (adp-viz) {
    \textbf{Adaptador Viz}\\[3pt]
    {\footnotesize Jupyter $\cdot$ \textsc{qgis}\\Transport-gis}
};
\node[adaptador, below=0.9cm of dom] (adp-per) {
    \textbf{Adaptador Persistencia}\\
    \texttt{src/infrastructure/}\\[3pt]
    {\footnotesize Redis $\cdot$ Celery\\PostgreSQL/PostGIS}
};

%% — Flechas entre adaptadores y dominio —
\draw[flecha] (adp-api) -- (dom);
\draw[flecha] (dom)     -- (adp-viz);
\draw[flecha] (dom)     -- (adp-per);

%% — Borde Docker (en fondo) —
\begin{scope}[on background layer]
    \node[draw=unamAzul!35, dashed, rounded corners=8pt, fill=unamAzul!3,
          fit={(adp-api)(adp-viz)(adp-per)(dom)},
          inner sep=16pt] (docker) {};
\end{scope}

%% — Etiqueta Docker centrada en la parte superior —
\node[font=\scriptsize\bfseries\color{unamAzul!65}, anchor=north]
    at (docker.north) {Docker};

\end{tikzpicture}
